\documentclass[aspectratio=169,11pt,xcolor=table,t]{beamer}
\usepackage{slidestyle}
\externaldocument[H-]{handout2}
\ifdefined\notesmode
  \setbeameroption{show only notes}
\else
  \setbeameroption{hide notes}
\fi
\newcommand\hp[1]{\fref{Handout p.\pageref{H-#1}}}

\begin{document}

% ================================================================ opening
\begin{frame}{Lesson 2 \textperiodcentered{} Exponentials and Logarithms}
\relax
NCEA Level 2 Mathematics \textperiodcentered{} Algebra\par\vspace{3mm}
{\bfseries By the end of today you can:}
\begin{enumerate}
\item switch between power form and log form
\item use the three log laws to combine and solve
\item solve equations with the unknown in the power
\item answer growth and decay questions in context
\end{enumerate}
\note{%
Teaching line, not a question slide.\par
100 分钟。时间：Do Now 0:00–0:05；核心思想 0:05–0:07；Block 1 log 的意思 0:07–0:17；Block 2 log laws 0:17–0:32；Block 3 解指数方程 0:32–0:47；Block 4 增长和衰减 0:47–0:57；Classwork 0:57–1:17；真题演练 1:17–1:35；收尾 1:35–1:40。\par
发四份纸：Handout 2、Classwork 2、Practice set 2、Homework 2。学生需要计算器（log、ln 键）。\par
先用 5 分钟对 Homework 1 的答案（按 L 编号说错在哪），再开始 Do Now；时间不够就把 Block 4 的 Practice 放进作业。\par
Say it:\par
exponential = ek-spoh-NEN-shul\par
logarithm = LOG-uh-rith-um\par
ln = el-EN}
\end{frame}

\begin{frame}{Do Now}
\relax
\qline{1} Simplify $\dfrac{(2x)^{3}}{x^{-1}}$.
\qline{2} Write $\sqrt[3]{x^{2}}$ as a power of $x$.
\qline{3} Solve $2^{x} = 32$.
\qline{4} Factorise $2x^{2} + 7x + 3$.
\qline{5} Which is bigger: $2^{10}$ or $10^{3}$?
\note{%
Q1 $8x^{4}$\par
Q2 $x^{\frac{2}{3}}$\par
Q3 $x = 5$\par
Q4 $(2x+1)(x+3)$\par
Q5 $2^{10} = 1024$ is bigger\par
Short solution:\par
Q1 $8x^{3} \div x^{-1} = 8x^{3-(-1)}$\par
Q3 2, 4, 8, 16, 32：数了 5 次\par
Watch for:\par
Q1 写 $8x^{2}$：减负数（L1）\par
Q4 $ac = 6$：6 和 1\par
Diagnostic:\par
Q1、Q2 错：Lesson 1 的指数还不稳，Block 2 讲 log laws 前先回顾 index laws\par
Q3 靠猜也可以；问他 “$2^{x} = 30$ 怎么办”，引出今天的 log\par
Q4 错：Block 3 的 disguised quadratic 要慢一点}
\end{frame}

\begin{frame}{The one idea}
\relax
\fbx{2^{3} = 8 \qquad\Longleftrightarrow\qquad \log_{2} 8 = 3}
A log answers one question: \textbf{what power?}\par\smallskip
$\log_{2} 8$ means: 2 to the power \gap[16mm] gives 8.
\ask{So what is $\log_{2} 16$?}
\hp{h2:meaning}
\note{%
Gap what (the power 3)\par
Q 4, because $2^{4} = 16$\par
Short solution:\par
Q log 的答案就是一个 power；把它读成问题：“2 的几次方是 16？”\par
Watch for:\par
答 8（$16 \div 2$）：把 log 当成除法\par
Ask:\par
$\log_{2} 2$？（1）$\log_{2} 1$？（0）}
\end{frame}

% ================================================================ block 1
\begin{frame}{Quick review: powers}
\relax
\qline{1} $2^{6} = \gap[14mm]$
\qline{2} $3^{4} = \gap[14mm]$
\qline{3} $10^{-2} = \gap[14mm]$
\qline{4} $5^{0} = \gap[14mm]$
\qline{5}
\fbx{9^{\frac{1}{2}} = \underline{\hspace{12mm}}}
\hp{h2:review}
\note{%
Q1 64\par
Q2 81\par
Q3 0.01\par
Q4 1\par
Q5 3\par
Short solution:\par
Q3 $10^{-2} = \dfrac{1}{10^{2}} = \dfrac{1}{100}$\par
Q5 power $\dfrac{1}{2}$ 是 square root\par
Watch for:\par
Q3 答 $-100$ 或 $-20$：负 index（L1）\par
讲义最后一页有常用的 power 表，今天每个 log 都会用到}
\end{frame}

\begin{frame}{Switching forms}
\relax
\fbx{y = b^{x} \qquad\Longleftrightarrow\qquad \log_{b} y = x}
Base to the bottom, $y$ in the brackets.\par\smallskip
\wline{3^{4} = 81 \quad\Longleftrightarrow\quad \log_{3} \underline{\hspace{9mm}} = \underline{\hspace{9mm}}}
\wline{\log_{5} x = 3 \quad\Longleftrightarrow\quad x = \underline{\hspace{14mm}}}
\ask{In $\log_{5} x = 3$, which number is the base?}
\hp{h2:meaning}
\note{%
Gap 81; 4; $5^{3} = 125$\par
Q 5\par
Short solution:\par
Q base 是 log 下面的小数字；它在 power form 里是底数\par
Watch for:\par
写成 $x = 3^{5}$：base 和 power 弄反（\Lref{logform}）\par
Say it:\par
base = BAYS}
\end{frame}

\begin{frame}{When the unknown is the base}
\relax
Solve $\log_{x} 49 = 2$.
\wline{x^{2} = 49}
\wline{x = 7 \quad\text{or}\quad x = -7\;?}
\ask{Can a base be negative? Which answer survives?}
\hp{h2:meaning}
\note{%
Q no; $x = 7$\par
Short solution:\par
Q base 必须是正数（且不是 1），所以 $x = -7$ 舍去，并写出理由\par
Watch for:\par
只写 $x = \pm 7$ 或只写 7 不说理由：评分表要 evidence of negative value being disregarded（\Lref{argument}）}
\end{frame}

\begin{frame}{Logs you can do in your head}
\relax
\wline{\log_{7} 1 = \underline{\hspace{9mm}} \qquad \log_{6} 6 = \underline{\hspace{9mm}}}
\wline{\log 1000 = \underline{\hspace{9mm}} \qquad \ln e = \underline{\hspace{9mm}}}
No base written: base 10 (the \textbf{log} key). $\ln$ is base $e \approx 2.718$ (the \textbf{ln} key).
\ask{Why is $\log_{7} 1$ the same as $\log_{6} 1$?}
\hp{h2:meaning}
\note{%
Gap 0; 1; 3; 1\par
Q any base to the power 0 is 1\par
Short solution:\par
Q $7^{0} = 1$，$6^{0} = 1$，所以都是 0\par
Watch for:\par
$\log 1000$ 答 100：把 log 当成除以 10\par
计算器上 log 是 base 10，ln 是 base $e$；今天的指数方程两个都会用}
\end{frame}

\begin{frame}{Practice}
\relax
\qline{Q1} Evaluate \ (a) $\log_{2} 32$ \quad (b) $\log_{9} 3$ \quad (c) $\log 0.01$
\qline{Q2} Solve $\log_{4} x = 3$.
\qline{Q3} Solve $\log_{x} 125 = 3$.
\hp{h2:meaning}
\note{%
Q1 (a) 5; (b) $\dfrac{1}{2}$; (c) $-2$\par
Q2 $x = 64$\par
Q3 $x = 5$\par
Short solution:\par
Q1 (b) 9 的几次方是 3？square root，所以 $\dfrac{1}{2}$ (u)\par
Q2 $x = 4^{3}$ (u)\par
Q3 $x^{3} = 125$，cube root (u)\par
Watch for:\par
Q1 (b) 答 3：方向反了（\Lref{logform}）\par
Q3 这里 cube root 只有一个值，不需要舍去}
\end{frame}

% ================================================================ block 2: laws
\begin{frame}{The three log laws}
\relax
\begin{rulebox}
$\log_{b}(xy) = \log_{b}x + \log_{b}y$\qquad (multiply: add the powers)\par
$\log_{b}\Bigl(\dfrac{x}{y}\Bigr) = \log_{b}x - \log_{b}y$\qquad (divide: subtract)\par
$\log_{b}(x^{n}) = n\log_{b}x$\qquad (power: multiply)
\end{rulebox}
\wline{\log_{2}3 + \log_{2}5 = \log_{2}\underline{\hspace{9mm}}\qquad \log_{2}12 - \log_{2}4 = \log_{2}\underline{\hspace{9mm}}}
\ask{Which index law is behind the first log law?}
\hp{h2:laws}
\note{%
Gap 15; 3\par
Q $x^{a}\times x^{b} = x^{a+b}$: multiplying adds the powers\par
Short solution:\par
Q log 就是 power，所以 log 的加法对应里面的乘法\par
三条都在 formula sheet 上，不用背，但要会用\par
Watch for:\par
$\log_{2}3 + \log_{2}5 = \log_{2}8$：加法对应乘法，不是加法（\Lref{loglaw}）}
\end{frame}

\begin{frame}{Writing a single log}
\relax
Write $2\log x + \log 5 - \log 10$ as a single log.
\wline{= \log x^{\underline{\hspace{7mm}}} + \log 5 - \log 10}
\wline{= \log\frac{5x^{2}}{10} = \log\underline{\hspace{14mm}}}
\ask{Why move the 2 up as a power first?}
\hp{h2:laws}
\note{%
Gap 2; $\dfrac{x^{2}}{2}$\par
Q the add and subtract laws only work on logs with nothing in front\par
Short solution:\par
Q 前面的系数先变成 power，然后加号进分子、减号进分母，最后约分\par
Watch for:\par
写成 $\log 2x$：2 没有变成 power（\Lref{bracketpower}）\par
答案停在 $\log\dfrac{5x^{2}}{10}$：分数要约到最简}
\end{frame}

\begin{frame}{Equations with two logs}
\relax
Solve $\log_{3}(x+4) + \log_{3}(x-4) = 2$.
\wline{\log_{3}\bigl((x+4)(x-4)\bigr) = 2}
\wline{x^{2} - 16 = 3^{2}}
\wline{x = 5 \quad\text{or}\quad x = -5}
\ask{Put $x = -5$ into $\log_{3}(x-4)$. What goes wrong?}
\hp{h2:laws}
\note{%
Q $\log_{3}(-9)$ does not exist: reject $x = -5$; $x = 5$\par
Short solution:\par
Q 先合并成一个 log，再换成 power form：$x^{2} - 16 = 9$\par
Watch for:\par
两个解都留下（\Lref{argument}）\par
右边的 2 换成 $3^{2} = 9$，不是 $2^{3}$（\Lref{logform}）}
\end{frame}

\begin{frame}{Spot the error (1)}
\relax
\textbf{Sam's answer has 2 errors.} Find each one, fix it, and say why it loses marks.
\begin{samcard}
\flawed{(a)\ \ Expand log(x + 3)\quad\ \ log(x + 3) = log x + log 3\par
(b)\ \ Write 3 log 2 as a single log\quad\ \ 3 log 2 = log 6}
\end{samcard}
\hp{h2:lose}
\note{%
1. (a) there is no law for $\log(x+3)$: it cannot be split (\Lref{loglaw})\par
2. (b) $3\log 2 = \log 2^{3} = \log 8$, not $\log 6$ (\Lref{bracketpower})\par
Short solution:\par
(a) $\log(x+3)$ already simplest\par
(b) $\log 8$\par
Watch for:\par
(a) 代 $x = 7$ 检查：$\log 10 = 1$，但 $\log 7 + \log 3 = \log 21 \approx 1.32$}
\end{frame}

\begin{frame}{Practice}
\relax
\qline{Q1} Write $\log(4a) + 3\log\Bigl(\dfrac{a}{2}\Bigr)$ as a single logarithm.
\qline{Q2} Solve $\log_{2}(x+2) + \log_{2}(x-2) = 5$.
\qline{Q3} Solve $\log x - \log(x-3) = \log 4$.
\hp{h2:laws}
\note{%
Q1 $\log\dfrac{a^{4}}{2}$\par
Q2 $x = 6$\par
Q3 $x = 4$\par
Short solution:\par
Q1 $\log 4a + \log\dfrac{a^{3}}{8} = \log\dfrac{4a^{4}}{8}$ (r)\par
Q2 $x^{2} - 4 = 32$，$x = \pm 6$，舍去 $-6$ (r)\par
Q3 $\log\dfrac{x}{x-3} = \log 4$，所以 $x = 4x - 12$ (u)\par
Watch for:\par
Q1 $\Bigl(\dfrac{a}{2}\Bigr)^{3}$ 分母也要立方（L2）\par
Q2 不舍去 $-6$（\Lref{argument}）}
\end{frame}

% ================================================================ block 3: solving
\begin{frame}{The unknown in the power}
\relax
Solve $3^{x} = 50$.
\wline{\log 3^{x} = \log 50}
\wline{x\log 3 = \log 50}
\wline{x = \frac{\log 50}{\log 3} = \underline{\hspace{16mm}}}
\ask{Why can't you just divide 50 by 3?}
\hp{h2:solve}
\note{%
Gap 3.561\par
Q $3^{x}$ is not $3 \times x$; the power has to come down first\par
Short solution:\par
Q 两边取 log，power law 把 $x$ 拿下来，再除以 $\log 3$\par
检查：$3^{3} = 27$，$3^{4} = 81$，所以答案在 3 和 4 之间\par
Watch for:\par
$\dfrac{\log 50}{\log 3}$ 写成 $\log\dfrac{50}{3}$（\Lref{loglaw}）}
\end{frame}

\begin{frame}{Bracket the power}
\relax
Solve $5^{2x-1} = 200$.
\wline{(2x-1)\log 5 = \log 200}
\wline{2x - 1 = \frac{\log 200}{\log 5} = 3.2920\ldots}
\wline{x = \underline{\hspace{16mm}}}
\ask{What goes wrong if you write $2x - 1\log 5$?}
\hp{h2:solve}
\note{%
Gap 2.146\par
Q only the 1 gets multiplied by $\log 5$\par
Short solution:\par
Q 整个 power $(2x-1)$ 拿下来要加括号；先求 $2x - 1$，最后再解 $x$\par
中间值保留计算器里的全部位数，最后才四舍五入（\Lref{context}）\par
Watch for:\par
没有括号（\Lref{bracketpower}）}
\end{frame}

\begin{frame}{Same base: solve exactly}
\relax
Solve $4^{x+1} = 8^{x}$.
\wline{4 = 2^{2} \quad\text{and}\quad 8 = 2^{\underline{\hspace{7mm}}}}
\wline{2^{2(x+1)} = 2^{3x}}
\wline{2x + 2 = 3x}
\ask{What base do 4 and 8 share? What is $x$?}
\hp{h2:solve}
\note{%
Gap 3\par
Q base 2; $x = 2$\par
Short solution:\par
Q 底数相同时 power 相等；精确解，不用计算器\par
检查：$4^{3} = 64 = 8^{2}$\par
Watch for:\par
$(2^{2})^{x+1}$ 写成 $2^{2x+1}$：括号（L1 的 power of a power）}
\end{frame}

\begin{frame}{$e$ and $\ln$}
\relax
Solve $e^{0.3t} = 5$.
\wline{\ln e^{0.3t} = \ln 5}
\wline{0.3t = \ln 5}
\wline{t = \underline{\hspace{16mm}}}
\ask{What undoes $e$ to a power?}
\hp{h2:solve}
\note{%
Gap 5.365\par
Q $\ln$, because $\ln e^{k} = k$\par
Short solution:\par
Q $t = \dfrac{\ln 5}{0.3}$；也可以用 log：$t = \dfrac{\log 5}{0.3 \log e}$，答案一样\par
Watch for:\par
写 $t = \ln\dfrac{5}{0.3}$：先除再取 ln 是错的\par
考试里常见 $T = T_{0} + ke^{-rt}$ 的冷却模型，先把 $e$ 的那一项单独移出来}
\end{frame}

\begin{frame}{A quadratic in disguise}
\relax
Solve $3^{2x} - 4\times 3^{x} + 3 = 0$.\par
Let $u = 3^{x}$. Then $3^{2x} = (3^{x})^{2} = u^{2}$.
\wline{u^{2} - 4u + 3 = 0}
\wline{(u-1)(u-3) = 0}
\wline{3^{x} = 1 \quad\text{or}\quad 3^{x} = 3}
\ask{What are the two values of $x$?}
\hp{h2:solve}
\note{%
Q $x = 0$ or $x = 1$\par
Short solution:\par
Q $3^{x} = 1$ 时 $x = 0$；$3^{x} = 3$ 时 $x = 1$ (t)\par
如果某个 $u$ 是负数，要舍去：$3^{x}$ 永远是正的\par
Watch for:\par
解出 $u$ 就停：要换回 $x$\par
认出 $3^{2x} = (3^{x})^{2}$ 是 Excellence 的关键一步}
\end{frame}

\begin{frame}{Spot the error (2)}
\relax
Solve $5^{x-2} = 40$.\par
\textbf{Sam's answer has 2 errors.} Find each one, fix it, and say why it loses marks.
\begin{samcard}
\flawed{x \hminus 2 log 5 = log 40\par
x \hminus 2 = \hfr{log 40}{log 5} = log 8 = 0.903\par
x = 2.903}
\end{samcard}
\hp{h2:lose}
\note{%
1. the power comes down in a bracket: $(x-2)\log 5 = \log 40$ (\Lref{bracketpower})\par
2. $\dfrac{\log 40}{\log 5} \neq \log 8$: divide on the calculator, $= 2.292$ (\Lref{loglaw})\par
Short solution:\par
$x - 2 = 2.292$，$x = 4.292$\par
Watch for:\par
检查：$5^{4.292-2} = 5^{2.292} \approx 40$；Sam 的 $5^{0.903} \approx 4.3$}
\end{frame}

\begin{frame}{Practice}
\relax
\qline{Q1} Solve $2^{x} = 300$.
\qline{Q2} Solve $7^{3x+1} = 90$.
\qline{Q3} Solve $9^{x} = 27^{x-1}$ exactly.
\hp{h2:solve}
\note{%
Q1 $x = 8.229$\par
Q2 $x = 0.4375$\par
Q3 $x = 3$\par
Short solution:\par
Q1 $x = \dfrac{\log 300}{\log 2}$ (u)\par
Q2 $3x + 1 = \dfrac{\log 90}{\log 7} = 2.3124$ (u)\par
Q3 $3^{2x} = 3^{3x-3}$，$2x = 3x - 3$ (r)\par
Watch for:\par
Q2 括号（\Lref{bracketpower}）\par
Q3 用计算器得到小数也对，但题目要 exactly}
\end{frame}

% ================================================================ block 4: growth
\begin{frame}{Growth and decay}
\relax
\fbx{P = A\,r^{t}}
A car bought for \$18\,000 loses 15\% of its value each year.
\wline{V = 18\,000 \times \underline{\hspace{12mm}}^{\,t}}
\ask{Why is the \Tgrowth{} 0.85 and not 0.15?}
\hp{h2:growth}
\note{%
Gap 0.85\par
Q each year it keeps 85\% of its value\par
Short solution:\par
Q 增长 $R\%$：乘 $1 + \dfrac{R}{100}$；减少 $R\%$：乘 $1 - \dfrac{R}{100}$\par
Watch for:\par
用 0.15：一年后只剩 15\%（\Lref{factor}）}
\end{frame}

\begin{frame}{When does it happen?}
\relax
When is the car first worth less than \$6000?
\wline{18\,000(0.85)^{t} = 6000}
\wline{0.85^{t} = \frac{1}{3}}
\wline{t = \frac{\log\frac{1}{3}}{\log 0.85} = 6.76}
\ask{Is the answer ``6.76''? What do you write?}
\hp{h2:growth}
\note{%
Q no: it is first worth less than \$6000 during the 7th year, after about 6.8 years\par
Short solution:\par
Q 先除以 18000，再取 log；最后一句话要有单位和情境\par
Watch for:\par
只写 6.76 或 “6 years”：没解释（\Lref{context}）\par
$\log\dfrac{1}{3}$ 和 $\log 0.85$ 都是负数，相除是正数}
\end{frame}

\begin{frame}{What rate?}
\relax
\$5000 grows to \$7500 in 8 years at a fixed annual rate $R\%$.
\wline{5000\Bigl(1 + \frac{R}{100}\Bigr)^{8} = 7500}
\fbx{1 + \frac{R}{100} = 1.5^{\frac{1}{8}} = \underline{\hspace{18mm}}}
\ask{When the unknown is the base, do you take logs or a root?}
\hp{h2:growth}
\note{%
Gap 1.0520\par
Q a root: the 8th root, then $R = 5.2\%$\par
Short solution:\par
Q 未知数在 base，不在 power：开 8 次方（power $\dfrac{1}{8}$）\par
Watch for:\par
答 $R = 1.052$ 或 0.052：要换成百分数\par
答 $R = 6.25$（把 50\% 平均分给 8 年）：复利不是平均}
\end{frame}

\begin{frame}{Spot the error (3)}
\relax
A \$15\,000 car loses 20\% of its value each year. When is it first worth less than \$5000?\par
\textbf{Sam's answer has 2 errors.} Find each one, fix it, and say why it loses marks.
\begin{samcard}
\flawed{15000(0.2)\hsp{t} = 5000\qquad 0.2\hsp{t} = \hfr{1}{3}\qquad t = \hfr{log(1/3)}{log 0.2} = 0.683\par
Answer: 0.683}
\end{samcard}
\hp{h2:lose}
\note{%
1. a 20\% loss keeps 80\%: the factor is 0.8, not 0.2 (\Lref{factor})\par
2. the answer needs units and a sentence in context (\Lref{context})\par
Short solution:\par
$0.8^{t} = \dfrac{1}{3}$，$t = 4.92$：during the 5th year, after about 4.9 years\par
Watch for:\par
Sam 的 0.683 年明显不合理：一年后车还值 \$12000}
\end{frame}

\begin{frame}{Practice}
\relax
\qline{Q1} A bacteria count is $N = 200(1.35)^{t}$ after $t$ hours. When does it first pass 10\,000?
\qline{Q2} A coffee cools as $T = 20 + 65e^{-0.08t}$ ($t$ in minutes). When is it 50\,\textdegree C?
\qline{Q3} A town grows from 4200 to 5600 people in 6 years at a constant rate. Find the annual rate.
\hp{h2:growth}
\note{%
Q1 after about 13.0 hours\par
Q2 after about 9.7 minutes\par
Q3 about 4.9\% per year\par
Short solution:\par
Q1 $1.35^{t} = 50$，$t = \dfrac{\log 50}{\log 1.35} = 13.04$ (u)\par
Q2 $65e^{-0.08t} = 30$，$-0.08t = \ln\dfrac{30}{65}$，$t = 9.665$ (r)\par
Q3 $r^{6} = \dfrac{4}{3}$，$r = 1.0491$ (r)\par
Watch for:\par
Q2 先减 20 再除 65：顺序反了得不到 $e$ 单独一项\par
每题最后写一句话（\Lref{context}）}
\end{frame}

\begin{frame}{Classwork}
\relax
Classwork 2 \textperiodcentered{} about 20 minutes
\begin{itemize}
\item \textbf{Part A}: switching forms, single logs, solving
\item \textbf{Part B}: two-log equations and a decay model
\item \textbf{Part C}: a quadratic in disguise, explained
\item \textbf{Extension}, if you finish
\end{itemize}
\note{%
Teaching line, not a question slide.\par
答案在 classwork2\_answers.pdf（同版面，红色答案）。\par
A1 5；A2 36；A3 $\log 10 = 1$；A4 $x = 2.410$；A5 $x = 2$\par
B1 $x = 2.548$；B2 $x = 4$（舍去 $-2$）；B3 during the 7th year\par
C：$x = 0$ or $x = 2$\par
Extension：$x = 2.710$\par
学生做的时候不要讲；收卷后先看 Part C 有没有写 $u = 2^{x}$ 和舍去的理由}
\end{frame}

% ================================================================ exam run
\begin{frame}{Practice}
\relax
\textbf{(c)} Jessica is investigating a compounding investment. She wants to know how long it would take for an investment of \$1000 to double in value to \$2000. She forms the following equation:
\fbx{2000 = 1000\Bigl(1 + \frac{R}{100}\Bigr)^{D}}
where $R$ is the rate of return on the investment, as a percentage, and $D$ is the time that the investment would take to double in value, in years.\par\smallskip
\textbf{(i)} If an investment takes 11 years to double in value, what is its rate of return?
\fref{Practice set p.1}
\note{%
(i) $R = 6.5\%$\par
Short solution:\par
(i) $\Bigl(1 + \dfrac{R}{100}\Bigr)^{11} = 2$，$1 + \dfrac{R}{100} = 2^{\frac{1}{11}} = 1.065$\par
评分：列出方程并用 11 次方根是 u，正确答案是 r\par
Watch for:\par
未知数在 base：开方，不是取 log}
\end{frame}

\begin{frame}{Practice}
\relax
\textbf{(ii)} By making $D$ the subject of the expression:
\fbx{2000 = 1000\Bigl(1 + \frac{R}{100}\Bigr)^{D}}
show that
\fbx{D = \frac{\log(2)}{\log\bigl(1 + \frac{R}{100}\bigr)}}
\fref{Practice set p.1}
\note{%
(ii) divide by 1000, take logs of both sides, bring $D$ down, divide\par
Short solution:\par
(ii) $\Bigl(1 + \dfrac{R}{100}\Bigr)^{D} = 2$\par
$\log\Bigl(1 + \dfrac{R}{100}\Bigr)^{D} = \log 2$\par
$D\log\Bigl(1 + \dfrac{R}{100}\Bigr) = \log 2$\par
评分：写出两边取 log 的一行（第 2 或第 3 行）是 u；完整推出给定式子是 r\par
Watch for:\par
Show that 题只写最后一行：没有分}
\end{frame}

\begin{frame}{Practice}
\relax
In her research, Jessica comes across a simple but \textbf{approximate} rule for calculating $D$, the time that the investment would take to double in value. It is commonly called the `Rule of 72', and it states that: $D = \dfrac{72}{R}$\par\smallskip
Jessica wonders how close the values of $D$ from the `Rule of 72' are to those calculated using the actual expression, which is:
\fbx{D = \frac{\log(2)}{\log\bigl(1 + \frac{R}{100}\bigr)}}
\fref{Practice set p.2}
\note{%
Teaching line, not a question slide.\par
这是 (iii) 的背景，先读题，下一页是问题\par
问学生：$R = 8$ 时 Rule of 72 给几年？（9 年）真实值？（$\dfrac{\log 2}{\log 1.08} = 9.006$ 年）\par
Say it:\par
approximate = uh-PROK-sih-mit}
\end{frame}

\begin{frame}{Practice}
\relax
\textbf{(iii)} Show clearly that the value of $R$ for which the `Rule of 72' exactly calculates $D$, is the solution to the equation:
\fbx{2^{R} - \Bigl(1 + \frac{R}{100}\Bigr)^{72} = 0}
You do not need to solve this equation.
\fref{Practice set p.2}
\note{%
(iii) set the two expressions for $D$ equal and rearrange with log laws\par
Short solution:\par
(iii) $\dfrac{72}{R} = \dfrac{\log 2}{\log\bigl(1 + \frac{R}{100}\bigr)}$\par
$72\log\Bigl(1 + \dfrac{R}{100}\Bigr) = R\log 2$\par
$\log\Bigl(1 + \dfrac{R}{100}\Bigr)^{72} = \log 2^{R}$，所以 $\Bigl(1 + \dfrac{R}{100}\Bigr)^{72} = 2^{R}$\par
评分：两式相等是 u；处理 power（第 3 行）是 r；推出给定式子且数学表述正确是 t\par
Watch for:\par
两边的 log 去掉之前要先写成 $\log(\ldots) = \log(\ldots)$ 的形式}
\end{frame}

\begin{frame}{Practice}
\relax
\qline{(a)} (ii)\quad Solve the equation $2^{x} = 2022$.
\fref{Practice set p.3}
\note{%
(ii) $x = 10.98$\par
Short solution:\par
(ii) $x\log 2 = \log 2022$；评分：Correct solution (u)；写 $\log_{2} 2022$ 也接受\par
Watch for:\par
检查：$2^{11} = 2048$，所以略小于 11}
\end{frame}

\begin{frame}{Practice}
\relax
\qline{(b)} Simplify the following expression fully, writing your answer as a single logarithm.
\fbx{\log(3a) + 2\log\Bigl(\frac{a}{6}\Bigr)}
\fref{Practice set p.3}
\note{%
(b) $\log\dfrac{a^{3}}{12}$\par
Short solution:\par
(b) $\log 3a + \log\dfrac{a^{2}}{36} = \log\dfrac{3a^{3}}{36}$\par
评分：分数没有化简是 u；正确的 single log 是 r\par
Watch for:\par
$\Bigl(\dfrac{a}{6}\Bigr)^{2}$ 分母是 36（L2）}
\end{frame}

\begin{frame}{Practice}
\relax
\textbf{(c)} Consider the equation $\log_{2}(x-a) - \log_{2}(x+a) = c$, where $a$ and $c$ are constants.\par\smallskip
\textbf{(i)} Show that when $x$ is made the subject of this equation,
\fbx{x = a\,\frac{1 + 2^{c}}{1 - 2^{c}}}
Ensure that you use correct mathematical statements in your reasoning.
\fref{Practice set p.4}
\note{%
(i) combine the logs, switch to power form, collect the $x$ terms, factorise\par
Short solution:\par
(i) $\log_{2}\dfrac{x-a}{x+a} = c$，所以 $\dfrac{x-a}{x+a} = 2^{c}$\par
$x - a = 2^{c}x + 2^{c}a$，$x(1 - 2^{c}) = a(1 + 2^{c})$\par
评分：合并 log 是 u；写出 power form（第 3 行）是 r；完整正确的推理是 t\par
Watch for:\par
这是 Lesson 1 的“换主元”：$x$ 在两边，要提公因式（L10）}
\end{frame}

\begin{frame}{Practice}
\relax
\textbf{(ii)} The equation $\log_{2}(x-a) - \log_{2}(x+a) = c$ is only possible to solve for some values of $a$ and for some values of $c$.\par\smallskip
Explaining your reasoning clearly, describe which values of $a$ and of $c$ will make the equation possible to solve.\par\smallskip
You may find it useful to recall that, when $x$ is made the subject of this equation,
\fbx{x = a\,\frac{1 + 2^{c}}{1 - 2^{c}}}
\fref{Practice set p.5}
\note{%
(ii) $c \neq 0$ and $a \neq 0$, with $x - a > 0$ and $x + a > 0$; this happens exactly when $a$ and $c$ have opposite signs\par
Short solution:\par
(ii) $c = 0$ 时 $1 - 2^{c} = 0$，$x$ 没有定义\par
$a = 0$ 时 $x = 0$，$\log_{2}0$ 没有定义\par
两个 log 的括号都要为正\par
评分：一个条件加理由是 r；两个条件加理由是 t\par
Watch for:\par
只写条件不写理由：Explaining your reasoning clearly\par
进一步：$a > 0$ 需要 $c < 0$；$a < 0$ 需要 $c > 0$（评分表只要求前面的条件）}
\end{frame}

% ================================================================ close
\begin{frame}{Ways to lose marks}
\relax
\begin{itemize}
\item \textbf{\Lref{logform}} Base to the bottom, $y$ in the brackets
\item \textbf{\Lref{loglaw}} No law for $\log(x+3)$; $\dfrac{\log 40}{\log 5}$ is not $\log 8$
\item \textbf{\Lref{bracketpower}} $(x-2)\log 5$: the whole power in a bracket
\item \textbf{\Lref{argument}} Check every answer: logs need positive brackets
\item \textbf{\Lref{factor}} A 20\% loss multiplies by 0.8
\item \textbf{\Lref{context}} A sentence with units; round at the end
\end{itemize}
\hp{h2:lose}
\note{%
Teaching line, not a question slide.\par
讲义上有 L11–L16 的完整说明\par
让学生对照 Classwork 和 Practice set 的错，说出是哪一个 L}
\end{frame}

\begin{frame}{Summary}
\relax
\begin{itemize}
\item A log is a power: $\log_{b} y = x$ means $b^{x} = y$.
\item Laws: plus multiplies, minus divides, a number in front becomes a power.
\item Unknown in the power: take logs, bracket the power.
\item Growth: $P = Ar^{t}$; answer in a sentence with units.
\end{itemize}
\hp{h2:meaning}
\note{%
Teaching line, not a question slide.\par
让学生用自己的话说第一条：“log 问的是什么？”}
\end{frame}

\begin{frame}{Homework}
\relax
Homework 2 \textperiodcentered{} about 60 minutes
\begin{itemize}
\item \textbf{Quick review}: four short items
\item \textbf{Question One}: meaning of logs and the log laws
\item \textbf{Question Two}: exponential equations
\item \textbf{Question Three}: growth and decay in context
\end{itemize}
Each question is marked like the real paper: N0 to E8.
\note{%
Teaching line, not a question slide.\par
答案在 hw2\_answers.pdf；每题的 (e) 是 Excellence 题\par
下节课开头 5 分钟对答案}
\end{frame}

\end{document}
